\begin{lexlean}{Quality}
\useglossary{lexlean.std.bool@1.2.0}
\useglossary{lexlean.std.nat@1.2.0}
\useglossary{lexleanmcp.titles@1.0.0}
\useglossary{prism.arch@1.0.0}
\useglossary{prism.qual@1.0.0}
\importmodule{Foundation.Qual}
\importmodule{Foundation.Arch}
\importmodule{Architecture}
\title{LexLean MCP quality model}
\begin{semanticmodule}
inductive{
  name: EvidenceStatus,
  type_parameters: [],
  parameters: [],
  constructors: [{name: Executed, fields: [String]}, {name: Pending, fields: [String]}]
}

structure{
  name: Named,
  type_parameters: [],
  parameters: [],
  fields: [{name: index, type: Nat}, {name: name, type: String}, {name: text, type: String}]
}

structure{
  name: Description,
  type_parameters: [],
  parameters: [],
  fields: [
    {name: index, type: Nat},
    {name: observation, type: String},
    {name: threshold, type: String}
  ]
}

structure{
  name: Cardinality,
  type_parameters: [],
  parameters: [],
  fields: [{name: measure, type: Nat}, {name: count, type: Nat}, {name: unit, type: String}]
}

structure{
  name: Realization,
  type_parameters: [],
  parameters: [],
  fields: [{name: requirement, type: Nat}, {name: component, type: Nat}]
}

structure{
  name: MeasureEvidence,
  type_parameters: [],
  parameters: [],
  fields: [{name: measure, type: Nat}, {name: status, type: EvidenceStatus}]
}

definition{
  name: characteristicFunctionalsuitability,
  parameters: [],
  result: Foundation.Qual::ProductQualityCharacteristic,
  body: (record Foundation.Qual::ProductQualityCharacteristic {index := 0})
}

definition{
  name: characteristicReliability,
  parameters: [],
  result: Foundation.Qual::ProductQualityCharacteristic,
  body: (record Foundation.Qual::ProductQualityCharacteristic {index := 1})
}

definition{
  name: characteristicPerformanceefficiency,
  parameters: [],
  result: Foundation.Qual::ProductQualityCharacteristic,
  body: (record Foundation.Qual::ProductQualityCharacteristic {index := 2})
}

definition{
  name: characteristicCompatibility,
  parameters: [],
  result: Foundation.Qual::ProductQualityCharacteristic,
  body: (record Foundation.Qual::ProductQualityCharacteristic {index := 3})
}

definition{
  name: characteristicMaintainability,
  parameters: [],
  result: Foundation.Qual::ProductQualityCharacteristic,
  body: (record Foundation.Qual::ProductQualityCharacteristic {index := 4})
}

definition{
  name: characteristicFlexibility,
  parameters: [],
  result: Foundation.Qual::ProductQualityCharacteristic,
  body: (record Foundation.Qual::ProductQualityCharacteristic {index := 5})
}

definition{
  name: characteristicInteractioncapability,
  parameters: [],
  result: Foundation.Qual::ProductQualityCharacteristic,
  body: (record Foundation.Qual::ProductQualityCharacteristic {index := 6})
}

definition{
  name: subcharacteristicFunctionalCorrectness,
  parameters: [],
  result: Foundation.Qual::QualitySubcharacteristic,
  body: (record Foundation.Qual::QualitySubcharacteristic {index := 0, characteristicIndex := 0})
}

definition{
  name: subcharacteristicFunctionalCompleteness,
  parameters: [],
  result: Foundation.Qual::QualitySubcharacteristic,
  body: (record Foundation.Qual::QualitySubcharacteristic {index := 1, characteristicIndex := 0})
}

definition{
  name: subcharacteristicFaultTolerance,
  parameters: [],
  result: Foundation.Qual::QualitySubcharacteristic,
  body: (record Foundation.Qual::QualitySubcharacteristic {index := 2, characteristicIndex := 1})
}

definition{
  name: subcharacteristicResourceUtilization,
  parameters: [],
  result: Foundation.Qual::QualitySubcharacteristic,
  body: (record Foundation.Qual::QualitySubcharacteristic {index := 3, characteristicIndex := 2})
}

definition{
  name: subcharacteristicInteroperability,
  parameters: [],
  result: Foundation.Qual::QualitySubcharacteristic,
  body: (record Foundation.Qual::QualitySubcharacteristic {index := 4, characteristicIndex := 3})
}

definition{
  name: subcharacteristicModularity,
  parameters: [],
  result: Foundation.Qual::QualitySubcharacteristic,
  body: (record Foundation.Qual::QualitySubcharacteristic {index := 5, characteristicIndex := 4})
}

definition{
  name: subcharacteristicInstallability,
  parameters: [],
  result: Foundation.Qual::QualitySubcharacteristic,
  body: (record Foundation.Qual::QualitySubcharacteristic {index := 6, characteristicIndex := 5})
}

definition{
  name: subcharacteristicLearnability,
  parameters: [],
  result: Foundation.Qual::QualitySubcharacteristic,
  body: (record Foundation.Qual::QualitySubcharacteristic {index := 7, characteristicIndex := 6})
}

definition{
  name: requirementFunctionalCorrectness,
  parameters: [],
  result: Foundation.Qual::QualityRequirement,
  body: (record Foundation.Qual::QualityRequirement {index := 0, subcharacteristicIndex := 0})
}

definition{
  name: requirementFunctionalCompleteness,
  parameters: [],
  result: Foundation.Qual::QualityRequirement,
  body: (record Foundation.Qual::QualityRequirement {index := 1, subcharacteristicIndex := 1})
}

definition{
  name: requirementFaultTolerance,
  parameters: [],
  result: Foundation.Qual::QualityRequirement,
  body: (record Foundation.Qual::QualityRequirement {index := 2, subcharacteristicIndex := 2})
}

definition{
  name: requirementResourceUtilization,
  parameters: [],
  result: Foundation.Qual::QualityRequirement,
  body: (record Foundation.Qual::QualityRequirement {index := 3, subcharacteristicIndex := 3})
}

definition{
  name: requirementInteroperability,
  parameters: [],
  result: Foundation.Qual::QualityRequirement,
  body: (record Foundation.Qual::QualityRequirement {index := 4, subcharacteristicIndex := 4})
}

definition{
  name: requirementModularity,
  parameters: [],
  result: Foundation.Qual::QualityRequirement,
  body: (record Foundation.Qual::QualityRequirement {index := 5, subcharacteristicIndex := 5})
}

definition{
  name: requirementInstallability,
  parameters: [],
  result: Foundation.Qual::QualityRequirement,
  body: (record Foundation.Qual::QualityRequirement {index := 6, subcharacteristicIndex := 6})
}

definition{
  name: requirementLearnability,
  parameters: [],
  result: Foundation.Qual::QualityRequirement,
  body: (record Foundation.Qual::QualityRequirement {index := 7, subcharacteristicIndex := 7})
}

definition{
  name: measureFunctionalCorrectness,
  parameters: [],
  result: Foundation.Qual::QualityMeasure,
  body: (record Foundation.Qual::QualityMeasure {index := 0, requirementIndex := 0, threshold := 0})
}

definition{
  name: measureFunctionalCompleteness,
  parameters: [],
  result: Foundation.Qual::QualityMeasure,
  body: (record Foundation.Qual::QualityMeasure {index := 1, requirementIndex := 1, threshold := 0})
}

definition{
  name: measureFaultTolerance,
  parameters: [],
  result: Foundation.Qual::QualityMeasure,
  body: (record Foundation.Qual::QualityMeasure {index := 2, requirementIndex := 2, threshold := 0})
}

definition{
  name: measureResourceUtilization,
  parameters: [],
  result: Foundation.Qual::QualityMeasure,
  body: (record Foundation.Qual::QualityMeasure {index := 3, requirementIndex := 3, threshold := 0})
}

definition{
  name: measureInteroperability,
  parameters: [],
  result: Foundation.Qual::QualityMeasure,
  body: (record Foundation.Qual::QualityMeasure {index := 4, requirementIndex := 4, threshold := 0})
}

definition{
  name: measureModularity,
  parameters: [],
  result: Foundation.Qual::QualityMeasure,
  body: (record Foundation.Qual::QualityMeasure {index := 5, requirementIndex := 5, threshold := 0})
}

definition{
  name: measureInstallability,
  parameters: [],
  result: Foundation.Qual::QualityMeasure,
  body: (record Foundation.Qual::QualityMeasure {index := 6, requirementIndex := 6, threshold := 0})
}

definition{
  name: measureLearnability,
  parameters: [],
  result: Foundation.Qual::QualityMeasure,
  body: (record Foundation.Qual::QualityMeasure {index := 7, requirementIndex := 7, threshold := 0})
}

definition{
  name: standardIso25010,
  parameters: [],
  result: Foundation.Arch::StandardReference,
  body: (record Foundation.Arch::StandardReference {standardIndex := 25010, editionIndex := 2023, standardId := "ISO/IEC 25010", edition := "2023"})
}

definition{
  name: characteristicRecords,
  parameters: [],
  result: (List Foundation.Qual::ProductQualityCharacteristic),
  body: (cons (characteristicFunctionalsuitability) (cons (characteristicReliability) (cons (characteristicPerformanceefficiency) (cons (characteristicCompatibility) (cons (characteristicMaintainability) (cons (characteristicFlexibility) (cons (characteristicInteractioncapability) (nil Foundation.Qual::ProductQualityCharacteristic))))))))
}

definition{
  name: subcharacteristicRecords,
  parameters: [],
  result: (List Foundation.Qual::QualitySubcharacteristic),
  body: (cons (subcharacteristicFunctionalCorrectness) (cons (subcharacteristicFunctionalCompleteness) (cons (subcharacteristicFaultTolerance) (cons (subcharacteristicResourceUtilization) (cons (subcharacteristicInteroperability) (cons (subcharacteristicModularity) (cons (subcharacteristicInstallability) (cons (subcharacteristicLearnability) (nil Foundation.Qual::QualitySubcharacteristic)))))))))
}

definition{
  name: requirementRecords,
  parameters: [],
  result: (List Foundation.Qual::QualityRequirement),
  body: (cons (requirementFunctionalCorrectness) (cons (requirementFunctionalCompleteness) (cons (requirementFaultTolerance) (cons (requirementResourceUtilization) (cons (requirementInteroperability) (cons (requirementModularity) (cons (requirementInstallability) (cons (requirementLearnability) (nil Foundation.Qual::QualityRequirement)))))))))
}

definition{
  name: measureRecords,
  parameters: [],
  result: (List Foundation.Qual::QualityMeasure),
  body: (cons (measureFunctionalCorrectness) (cons (measureFunctionalCompleteness) (cons (measureFaultTolerance) (cons (measureResourceUtilization) (cons (measureInteroperability) (cons (measureModularity) (cons (measureInstallability) (cons (measureLearnability) (nil Foundation.Qual::QualityMeasure)))))))))
}

definition{
  name: characteristicNames,
  parameters: [],
  result: (List Named),
  body: (cons (record Named {index := 0, name := "Functional suitability", text := "A product quality characteristic of ISO/IEC 25010:2023."}) (cons (record Named {index := 1, name := "Reliability", text := "A product quality characteristic of ISO/IEC 25010:2023."}) (cons (record Named {index := 2, name := "Performance efficiency", text := "A product quality characteristic of ISO/IEC 25010:2023."}) (cons (record Named {index := 3, name := "Compatibility", text := "A product quality characteristic of ISO/IEC 25010:2023."}) (cons (record Named {index := 4, name := "Maintainability", text := "A product quality characteristic of ISO/IEC 25010:2023."}) (cons (record Named {index := 5, name := "Flexibility", text := "A product quality characteristic of ISO/IEC 25010:2023."}) (cons (record Named {index := 6, name := "Interaction capability", text := "A product quality characteristic of ISO/IEC 25010:2023."}) (nil Named))))))))
}

definition{
  name: subcharacteristicNames,
  parameters: [],
  result: (List Named),
  body: (cons (record Named {index := 0, name := "Functional correctness", text := "A subcharacteristic of Functional suitability."}) (cons (record Named {index := 1, name := "Functional completeness", text := "A subcharacteristic of Functional suitability."}) (cons (record Named {index := 2, name := "Fault tolerance", text := "A subcharacteristic of Reliability."}) (cons (record Named {index := 3, name := "Resource utilization", text := "A subcharacteristic of Performance efficiency."}) (cons (record Named {index := 4, name := "Interoperability", text := "A subcharacteristic of Compatibility."}) (cons (record Named {index := 5, name := "Modularity", text := "A subcharacteristic of Maintainability."}) (cons (record Named {index := 6, name := "Installability", text := "A subcharacteristic of Flexibility."}) (cons (record Named {index := 7, name := "Learnability", text := "A subcharacteristic of Interaction capability."}) (nil Named)))))))))
}

definition{
  name: requirementNames,
  parameters: [],
  result: (List Named),
  body: (cons (record Named {index := 0, name := "Q1", text := "a tool result equals the engine's result for the same project bytes"}) (cons (record Named {index := 1, name := "Q2", text := "each in-scope API function is reached by exactly one tool"}) (cons (record Named {index := 2, name := "Q3", text := "no verification record for a verification that did not end `verified`"}) (cons (record Named {index := 3, name := "Q4", text := "check, snapshot, format, build start no child process"}) (cons (record Named {index := 4, name := "Q5", text := "the MCP implementation conforms to 2026-07-28"}) (cons (record Named {index := 5, name := "Q6", text := "protocol and tool semantics live only in `.lex.tex`"}) (cons (record Named {index := 6, name := "Q7", text := "a draft yields the same `build_id` on every supported host (the build ID hashes only the source and semantic IDs, LexLean SPEC §21.5)"}) (cons (record Named {index := 7, name := "Q8", text := "every diagnostic code in a result resolves through F11"}) (nil Named)))))))))
}

definition{
  name: descriptions,
  parameters: [],
  result: (List Description),
  body: (cons (record Description {index := 0, observation := "canonical JSON of the structured content of `lexlean_check`, `lexlean_build`, `lexlean_verify` against the CLI's command result for the same command, over every example and the 273 negative fixtures of AC-07", threshold := "all equal"}) (cons (record Description {index := 1, observation := "map from `Engine::{check, snapshot, format, build, verify}` and `sdk::{assemble, write_module, read_module, declaration_spans}` to tools", threshold := "total"}) (cons (record Description {index := 2, observation := "fault cases: child killed, verifier killed, store full", threshold := "zero records"}) (cons (record Description {index := 3, observation := "child processes observed by host process accounting during each call", threshold := "zero"}) (cons (record Description {index := 4, observation := "`mcp-conformance-server-2026-07-28` and `mcp-schema-2026-07-28`", threshold := "50 of 50 scenarios pass; zero schema failures"}) (cons (record Description {index := 5, observation := "tool or method names in host adapter source", threshold := "zero"}) (cons (record Description {index := 6, observation := "`build_id` on `linux/amd64` against `linux/arm64`", threshold := "equal"}) (cons (record Description {index := 7, observation := "codes resolved against codes returned", threshold := "all"}) (nil Description)))))))))
}

definition{
  name: cardinalities,
  parameters: [],
  result: (List Cardinality),
  body: (cons (record Cardinality {measure := 0, count := 273, unit := "negative fixtures of AC-07"}) (cons (record Cardinality {measure := 1, count := 9, unit := "in-scope API functions of the pinned LexLean"}) (cons (record Cardinality {measure := 4, count := 50, unit := "scenarios"}) (cons (record Cardinality {measure := 7, count := 86, unit := "diagnostic codes of the pinned LexLean registry"}) (nil Cardinality)))))
}

definition{
  name: realizations,
  parameters: [],
  result: (List Realization),
  body: (cons (record Realization {requirement := 0, component := 3}) (cons (record Realization {requirement := 0, component := 4}) (cons (record Realization {requirement := 1, component := 3}) (cons (record Realization {requirement := 2, component := 5}) (cons (record Realization {requirement := 2, component := 6}) (cons (record Realization {requirement := 3, component := 4}) (cons (record Realization {requirement := 4, component := 1}) (cons (record Realization {requirement := 4, component := 2}) (cons (record Realization {requirement := 5, component := 1}) (cons (record Realization {requirement := 6, component := 4}) (cons (record Realization {requirement := 7, component := 3}) (nil Realization))))))))))))
}

definition{
  name: evidence,
  parameters: [],
  result: (List MeasureEvidence),
  body: (cons (record MeasureEvidence {measure := 0, status := (ctor EvidenceStatus.Pending "Run where the tools are built: #10 for check and build, #11 for verify.")}) (cons (record MeasureEvidence {measure := 1, status := (ctor EvidenceStatus.Executed "Executed by the API-to-tool map as built: QM-01, QM-02 and QM-03 recompute the 33 public functions of the pinned LexLean surface from its vendored sources and require each to reach one tool the issues define or to state why no tool does; that each mapped tool exists and is reached is in pendingClauses (#9 for the sdk functions, #10 and #11 for the Engine calls).")}) (cons (record MeasureEvidence {measure := 2, status := (ctor EvidenceStatus.Pending "Run where verification runs as tasks (#11); the store-full case needs the store (#13).")}) (cons (record MeasureEvidence {measure := 3, status := (ctor EvidenceStatus.Pending "Run where the tools are built (#10).")}) (cons (record MeasureEvidence {measure := 4, status := (ctor EvidenceStatus.Pending "Run against the deployed core (#17), through the MCP oracles.")}) (cons (record MeasureEvidence {measure := 5, status := (ctor EvidenceStatus.Pending "Run once the host adapter is generated (#7).")}) (cons (record MeasureEvidence {measure := 6, status := (ctor EvidenceStatus.Pending "Run where build is a tool (#10), on both platforms.")}) (cons (record MeasureEvidence {measure := 7, status := (ctor EvidenceStatus.Executed "Executed by the explained codes as built: QM-04, QM-05 and QM-06 require each of the 86 codes of the pinned LexLean registry to have one complete entry that lexlean explain prints, and every code the negative fixtures return to resolve to one; that a code in a result of the running tool resolves through the running lexlean_explain is executed by TL-13, which runs where the tools are built (gateway/tool-tests), so no clause of Q8 is pending.")}) (nil MeasureEvidence)))))))))
}

definition{
  name: pendingClauses,
  parameters: [],
  result: (List Named),
  body: (cons (record Named {index := 1, name := "Q2 tools exist", text := "each in-scope function is reached by a tool that exists and runs; built where the tools are built (#9 for the sdk functions, #10 and #11 for the Engine calls), so no register row executes it yet"}) (nil Named))
}

definition{
  name: nonEmpty,
  parameters: [text : String],
  result: Bool,
  body: (if (prim equal text "" : Bool) then false else true)
}

definition{
  name: namedIndexes,
  parameters: [records : (List Named)],
  result: (List Nat),
  body: (match records
    | List.nil => (nil Nat)
    | List.cons head rest => (cons (project head index) (namedIndexes rest))),
  recursive_argument: records
}

definition{
  name: characteristicIndexes,
  parameters: [records : (List Foundation.Qual::ProductQualityCharacteristic)],
  result: (List Nat),
  body: (match records
    | List.nil => (nil Nat)
    | List.cons head rest => (cons (project head index) (characteristicIndexes rest))),
  recursive_argument: records
}

definition{
  name: subcharacteristicIndexes,
  parameters: [records : (List Foundation.Qual::QualitySubcharacteristic)],
  result: (List Nat),
  body: (match records
    | List.nil => (nil Nat)
    | List.cons head rest => (cons (project head index) (subcharacteristicIndexes rest))),
  recursive_argument: records
}

definition{
  name: subcharacteristicParents,
  parameters: [records : (List Foundation.Qual::QualitySubcharacteristic)],
  result: (List Nat),
  body: (match records
    | List.nil => (nil Nat)
    | List.cons head rest => (cons (project head characteristicIndex) (subcharacteristicParents rest))),
  recursive_argument: records
}

definition{
  name: requirementIndexes,
  parameters: [records : (List Foundation.Qual::QualityRequirement)],
  result: (List Nat),
  body: (match records
    | List.nil => (nil Nat)
    | List.cons head rest => (cons (project head index) (requirementIndexes rest))),
  recursive_argument: records
}

definition{
  name: requirementParents,
  parameters: [records : (List Foundation.Qual::QualityRequirement)],
  result: (List Nat),
  body: (match records
    | List.nil => (nil Nat)
    | List.cons head rest => (cons (project head subcharacteristicIndex) (requirementParents rest))),
  recursive_argument: records
}

definition{
  name: measureIndexes,
  parameters: [records : (List Foundation.Qual::QualityMeasure)],
  result: (List Nat),
  body: (match records
    | List.nil => (nil Nat)
    | List.cons head rest => (cons (project head index) (measureIndexes rest))),
  recursive_argument: records
}

definition{
  name: measureRequirements,
  parameters: [records : (List Foundation.Qual::QualityMeasure)],
  result: (List Nat),
  body: (match records
    | List.nil => (nil Nat)
    | List.cons head rest => (cons (project head requirementIndex) (measureRequirements rest))),
  recursive_argument: records
}

definition{
  name: measureThresholds,
  parameters: [records : (List Foundation.Qual::QualityMeasure)],
  result: (List Nat),
  body: (match records
    | List.nil => (nil Nat)
    | List.cons head rest => (cons (project head threshold) (measureThresholds rest))),
  recursive_argument: records
}

definition{
  name: descriptionIndexes,
  parameters: [records : (List Description)],
  result: (List Nat),
  body: (match records
    | List.nil => (nil Nat)
    | List.cons head rest => (cons (project head index) (descriptionIndexes rest))),
  recursive_argument: records
}

definition{
  name: cardinalityMeasures,
  parameters: [records : (List Cardinality)],
  result: (List Nat),
  body: (match records
    | List.nil => (nil Nat)
    | List.cons head rest => (cons (project head measure) (cardinalityMeasures rest))),
  recursive_argument: records
}

definition{
  name: realizationRequirements,
  parameters: [records : (List Realization)],
  result: (List Nat),
  body: (match records
    | List.nil => (nil Nat)
    | List.cons head rest => (cons (project head requirement) (realizationRequirements rest))),
  recursive_argument: records
}

definition{
  name: realizationComponents,
  parameters: [records : (List Realization)],
  result: (List Nat),
  body: (match records
    | List.nil => (nil Nat)
    | List.cons head rest => (cons (project head component) (realizationComponents rest))),
  recursive_argument: records
}

definition{
  name: evidenceMeasures,
  parameters: [records : (List MeasureEvidence)],
  result: (List Nat),
  body: (match records
    | List.nil => (nil Nat)
    | List.cons head rest => (cons (project head measure) (evidenceMeasures rest))),
  recursive_argument: records
}

definition{
  name: descriptionsStated,
  parameters: [records : (List Description)],
  result: Bool,
  body: (match records
    | List.nil => true
    | List.cons head rest => (((nonEmpty (project head observation)) && (nonEmpty (project head threshold))) && (descriptionsStated rest))),
  recursive_argument: records
}

definition{
  name: allZero,
  parameters: [values : (List Nat)],
  result: Bool,
  body: (match values | List.nil => true | List.cons head rest => ((head =? 0) && (allZero rest))),
  recursive_argument: values
}

definition{
  name: countOf,
  parameters: [wanted : Nat, rows : (List Cardinality)],
  result: Nat,
  body: (match rows
    | List.nil => 0
    | List.cons head rest => (if ((project head measure) =? wanted) then (project head count) else (countOf wanted rest))),
  recursive_argument: rows
}

definition{
  name: evidenceHonest,
  parameters: [records : (List MeasureEvidence)],
  result: Bool,
  body: (match records
    | List.nil => true
    | List.cons head rest => ((match (project head status)
      | EvidenceStatus.Executed reference => (nonEmpty reference)
      | EvidenceStatus.Pending reason => (nonEmpty reason)) && (evidenceHonest rest))),
  recursive_argument: records
}

definition{
  name: pendingMeasures,
  parameters: [records : (List MeasureEvidence)],
  result: (List Nat),
  body: (match records
    | List.nil => (nil Nat)
    | List.cons head rest => (match (project head status)
      | EvidenceStatus.Executed reference => (pendingMeasures rest)
      | EvidenceStatus.Pending reason => (cons (project head measure) (pendingMeasures rest)))),
  recursive_argument: records
}

definition{
  name: executedMeasures,
  parameters: [records : (List MeasureEvidence)],
  result: (List Nat),
  body: (match records
    | List.nil => (nil Nat)
    | List.cons head rest => (match (project head status)
      | EvidenceStatus.Executed reference => (cons (project head measure) (executedMeasures rest))
      | EvidenceStatus.Pending reason => (executedMeasures rest))),
  recursive_argument: records
}

definition{
  name: viewComponents,
  parameters: [wanted : Nat, specs : (List Architecture::ViewSpec)],
  result: (List Nat),
  body: (match specs
    | List.nil => (nil Nat)
    | List.cons head rest => (if ((project head view) =? wanted) then (project head components) else (viewComponents wanted rest))),
  recursive_argument: specs
}

definition{
  name: qualityViewComponents,
  parameters: [],
  result: (List Nat),
  body: (viewComponents 2 (Architecture::viewSpecs))
}

definition{
  name: editionIs2023,
  parameters: [],
  result: Bool,
  body: (((project (standardIso25010) standardIndex) =? 25010) && ((project (standardIso25010) editionIndex) =? 2023))
}

definition{
  name: qualityChainHolds,
  parameters: [],
  result: Bool,
  body: ((Architecture::allMembers (subcharacteristicParents (subcharacteristicRecords)) (characteristicIndexes (characteristicRecords))) && ((Architecture::allMembers (characteristicIndexes (characteristicRecords)) (subcharacteristicParents (subcharacteristicRecords))) && ((Architecture::allMembers (requirementParents (requirementRecords)) (subcharacteristicIndexes (subcharacteristicRecords))) && ((Architecture::allMembers (subcharacteristicIndexes (subcharacteristicRecords)) (requirementParents (requirementRecords))) && ((Architecture::allMembers (measureRequirements (measureRecords)) (requirementIndexes (requirementRecords))) && ((Architecture::allMembers (requirementIndexes (requirementRecords)) (measureRequirements (measureRecords))) && (allZero (measureThresholds (measureRecords)))))))))
}

definition{
  name: qualityValid,
  parameters: [],
  result: Bool,
  body: ((qualityChainHolds) && ((Architecture::sameNats (namedIndexes (characteristicNames)) (characteristicIndexes (characteristicRecords))) && ((Architecture::sameNats (namedIndexes (subcharacteristicNames)) (subcharacteristicIndexes (subcharacteristicRecords))) && ((Architecture::sameNats (namedIndexes (requirementNames)) (requirementIndexes (requirementRecords))) && ((Architecture::sameNats (descriptionIndexes (descriptions)) (measureIndexes (measureRecords))) && ((descriptionsStated (descriptions)) && ((Architecture::allMembers (cardinalityMeasures (cardinalities)) (measureIndexes (measureRecords))) && ((Architecture::allMembers (requirementIndexes (requirementRecords)) (realizationRequirements (realizations))) && ((Architecture::allMembers (realizationRequirements (realizations)) (requirementIndexes (requirementRecords))) && ((Architecture::allMembers (realizationComponents (realizations)) (Architecture::componentIndexes (Architecture::componentRecords))) && ((Architecture::allMembers (realizationComponents (realizations)) (qualityViewComponents)) && ((Architecture::sameNats (evidenceMeasures (evidence)) (measureIndexes (measureRecords))) && ((evidenceHonest (evidence)) && ((Architecture::allMembers (namedIndexes (pendingClauses)) (executedMeasures (evidence))) && (editionIs2023)))))))))))))))
}

theorem{
  name: qualityModelHolds,
  parameters: [],
  statement: ((qualityValid) = true),
  proof: decide{}
}

theorem{
  name: qualityChainIsComplete,
  parameters: [],
  statement: ((qualityChainHolds) = true),
  proof: decide{}
}

theorem{
  name: everyRequirementHasMeasure,
  parameters: [],
  statement: ((Architecture::allMembers (requirementIndexes (requirementRecords)) (measureRequirements (measureRecords))) = true),
  proof: decide{}
}

theorem{
  name: everyThresholdIsAllOrNone,
  parameters: [],
  statement: ((allZero (measureThresholds (measureRecords))) = true),
  proof: decide{}
}

theorem{
  name: requirementsAreQ1ToQ8,
  parameters: [],
  statement: ((requirementIndexes (requirementRecords)) = (cons 0 (cons 1 (cons 2 (cons 3 (cons 4 (cons 5 (cons 6 (cons 7 (nil Nat)))))))))),
  proof: decide{}
}

theorem{
  name: fixturesRequired,
  parameters: [],
  statement: ((countOf 0 (cardinalities)) = 273),
  proof: decide{}
}

theorem{
  name: scenariosRequired,
  parameters: [],
  statement: ((countOf 4 (cardinalities)) = 50),
  proof: decide{}
}

theorem{
  name: measuresPendingExceptQ2AndQ8,
  parameters: [],
  statement: ((pendingMeasures (evidence)) = (cons 0 (cons 2 (cons 3 (cons 4 (cons 5 (cons 6 (nil Nat)))))))),
  proof: decide{}
}

theorem{
  name: onlyApiMapAndExplainAreExecuted,
  parameters: [],
  statement: ((executedMeasures (evidence)) = (cons 1 (cons 7 (nil Nat)))),
  proof: decide{}
}

theorem{
  name: onlyQ2KeepsAPendingClause,
  parameters: [],
  statement: ((namedIndexes (pendingClauses)) = (cons 1 (nil Nat))),
  proof: decide{}
}

theorem{
  name: inScopeFunctionsAreCounted,
  parameters: [],
  statement: ((countOf 1 (cardinalities)) = 9),
  proof: decide{}
}

theorem{
  name: diagnosticCodesAreCounted,
  parameters: [],
  statement: ((countOf 7 (cardinalities)) = 86),
  proof: decide{}
}

theorem{
  name: requirementsRealizedInQualityView,
  parameters: [],
  statement: ((Architecture::allMembers (realizationComponents (realizations)) (qualityViewComponents)) = true),
  proof: decide{}
}

theorem{
  name: editionIsIso25010Of2023,
  parameters: [],
  statement: ((editionIs2023) = true),
  proof: decide{}
}

theorem{
  name: sevenCharacteristics,
  parameters: [],
  statement: ((characteristicIndexes (characteristicRecords)) = (cons 0 (cons 1 (cons 2 (cons 3 (cons 4 (cons 5 (cons 6 (nil Nat))))))))),
  proof: decide{}
}
\end{semanticmodule}
\end{lexlean}
